 \section{Conclusion}
\label{sec:conclusion}

This paper studied energy-efficient three-dimensional speed scheduling for UAV data collection from sequentially ordered ground nodes.
Departing from prior formulations that either fix the flight altitude or treat altitude changes as discrete transitions between horizontal segments, we formulated the joint horizontal-vertical speed optimization (\textit{JHVSO}) problem, in which the UAV's horizontal and vertical velocity components are scheduled jointly and continuously, the ground-to-air communication range varies with the flight altitude, and each ground node imposes its own minimum data collection duration.

The solution of the \textit{JHVSO} problem was built on a sequence of analytical results.
Exploiting the strict convexity of the horizontal and vertical power components, we proved that constant-velocity straight-line flight is the unique energy minimizer between any two states under a fixed duration, and derived the pitch-angle-dependent speed that minimizes the energy consumption per unit distance when the flight duration is unconstrained.
We further established the non-aftereffect property of the multi-stage state transitions, which legitimizes a dynamic programming formulation.
Building on these results, the proposed Optimal Reached Energy (ORE) algorithm decomposes each inter-node transition into gap and task segments, evaluates feasible transitions through the communication-boundary entry states, and computes the global energy-optimal trajectory over the discretized state space; an exact algebraic restructuring of the recurrence reduces the per-stage complexity from $O(m \times j \times n)$ to $O(m \times j + j \times n)$ without sacrificing optimality. A sensitivity study further confirms that the discretized optimum converges to the continuous optimum with negligible loss (within $0.07\%$ at the operating resolution).

Extensive simulations under varying channel conditions, collection durations, node densities, and inter-node distances show that \textit{JHVSO} consistently consumes less energy than ASSUME, MSMO, ACO, and a greedy speed-height baseline, with an energy reduction of up to approximately $16.8\%$ over the greedy speed–height baseline and by up to 5.8\% over the strongest competing baseline.
A case study based on real-world traffic-camera locations further illustrates the practical trajectory patterns produced by ORE, in which the UAV cruises near its maximum-range speed along gap segments and adapts both its speed and altitude to the collection requirement and coverage geometry of each ground node.

Overall, by jointly scheduling the horizontal and vertical speed components together with the flight altitude, the proposed JHVSO framework delivers a rigorous, globally optimal, and computationally efficient solution to energy-minimal UAV data collection, reconciling the aerodynamic flight power model with the altitude-dependent communication coverage.


% Numbered list
% Use the style of numbering in square brackets.
% If nothing is used, default style will be taken.
%\begin{enumerate}[a)]
%\item 
%\item 
%\item 
%\end{enumerate}  

% Unnumbered list
%\begin{itemize}
%\item 
%\item 
%\item 
%\end{itemize}  

% Description list
%\begin{description}
%\item[]
%\item[] 
%\item[] 
%\end{description}  

% \clearpage %%Remove this from your manuscript

% Figure
% \begin{figure}%[]
%   \centering
% %    \includegraphics{}
%     \caption{}\label{fig1}
% \end{figure}


% \begin{table}%[]
% \caption{}\label{tbl1}
% \begin{tabular*}{\tblwidth}{@{}LL@{}}
% \toprule
%   &  \\ % Table header row
% \midrule
%  & \\
%  & \\
%  & \\
%  & \\
% \bottomrule
% \end{tabular*}
% \end{table}

% Uncomment and use as the case may be
%\begin{theorem} 
%\end{theorem}

% Uncomment and use as the case may be
%\begin{lemma} 
%\end{lemma}

%%%%%%%%%%%%%%%%%%%%%%%%%%%%%%%%%%%%%%%%%%%%%%%%%%%%%%%%%%%%%%%%%%%%%%%%%%%%%%%%%%
%%%%%%%%%%%%%%%%%%%%%%%%%%%%%%%%%%%%%%%%%%%%%%%%%%%%%%%%%%%%%%%%%%%%%%%%%%%%%%%%%%

%% The Appendices part is started with the command \appendix;
%% appendix sections are then done as normal sections
